\documentclass[10pt,a4paper]{amsart}
\usepackage[margin=19mm]{geometry}
\usepackage{amsmath,amssymb,mathtools}
\usepackage[scr=rsfs,cal=euler]{mathalfa}
\usepackage{xfrac}
\usepackage[unicode,hidelinks,bookmarks=false]{hyperref}
\newcommand{\ie}{i.e.\ }
\newcommand{\cf}{cf.\ }
\newcommand{\resp}{respectively\ }
\newcommand{\Subsection}{\subsection}
\providecommand{\nobreakdash}{\nobreak\hbox{-}}
\setlength{\emergencystretch}{2em}
\multlinegap0pt
\usepackage{bm}
\usepackage{bbm}
\usepackage{dsfont}
\usepackage{xspace}
\usepackage{cancel}
\usepackage{mathtools}
\usepackage{amsthm}
\newtheorem{theorem}{Theorem}
\newcommand{\Hom}{\operatorname{Hom}}
\title[Ideal Structures in Idealizations and Cardinalities of Annih]{Ideal Structures in Idealizations and Cardinalities of Annihilating Ideals (Theorem 2.1)}
\author{Reza Nikandish}
\date{Source: arXiv:2609.25157v1 (CC BY 4.0)}
\begin{document}
\maketitle

\noindent\emph{Source and licence.} Ideal Structures in Idealizations and Cardinalities of Annihilating Ideals. Version arXiv:2609.25157v1 (CC BY 4.0), distributed under the Creative Commons Attribution 4.0 International licence, as recorded in the arXiv licence metadata for this exact version and shown on the abstract page https://arxiv.org/abs/2609.25157v1. Full rights notice: licenses/arxiv-2609.25157-CC-BY-4.0.txt.\par\smallskip

\noindent\textit{Editorial scope.} Counted unit: the Theorem printed as Theorem 2.1 of the article (source0.tex lines 227--275, source label \texttt{thm:classification}), delivered together with the article's own complete original proof of it (source0.tex lines 277--410, one \texttt{proof} environment of 134 lines). No mathematics is written, repaired or re-derived: statement and proof are the article's own text line for line. Required context retained uncounted: none. Every definition, display and label of those lines is retained unchanged. Omitted from this package and not redistributed: 1--226, 276--276, 411--1096. Nothing inside a retained excerpt is omitted or altered: every retained body is delivered exactly as the article's lines are written, so no omission receipt of any kind accompanies this package. Citations: the retained excerpts carry no citation command at all, so no bibliography entry of the article is transcribed and no reference is redistributed. Reference adaptation: none was needed and none was made; every reference inside the delivered unit resolves inside it, so no unresolved reference or citation is produced. Printed slips kept literal: no printed word, symbol, hypothesis or number of the counted statement or of its proof was altered. Layout and class: the article's own class is \texttt{article} with packages including geometry, amsmath, amssymb, amsthm, mathtools, enumitem, microtype, hyperref; the wrapper is the lane's pinned \texttt{amsart} editorial wrapper at 10pt on A4, which restates the theorem-like environments the retained text needs and adds the article's own macro definitions that the retained text needs (\texttt{\textbackslash Hom}), with extra wrapper packages (bm, bbm, dsfont, xspace, cancel, mathtools). The delivered numbering is the unit's own. Assets: no retained span contains a figure environment, a graphics-inclusion command, a PDF-page inclusion, a table or a TikZ picture; no image file is copied into this package, the package declares no asset dependency and no image file is redistributed with it. Licence: version arXiv:2609.25157v1 of this article is distributed under the Creative Commons Attribution 4.0 International licence, as recorded in the arXiv licence metadata for this exact version and shown on the abstract page https://arxiv.org/abs/2609.25157v1. A full rights notice is delivered at licenses/arxiv-2609.25157-CC-BY-4.0.txt. Characters: every retained excerpt is pure ASCII, so the delivered engine, XeLaTeX with its default Latin Modern text font, reports no missing character.
	\begin{theorem}
		\label{thm:classification}
		Let $D$ be a commutative ring, let $M$ be a $D$-module, and put
		\[
		R=D\ltimes M.
		\]
		For an ideal $J$ of $R$, define
		\[
		I=\pi(J)
		\]
		where $\pi:R\to D$ is the natural projection, and put
		\[
		N=\{m\in M:(0,m)\in J\}.
		\]
		Then $I$ is an ideal of $D$, $N$ is a $D$-submodule of $M$, and
		\[
		IM\subseteq N.
		\]
		Moreover, there is a unique homomorphism
		\[
		\varphi\in\Hom_D(I,M/N)
		\]
		such that
		\[
		J
		=
		\{(d,m)\in D\ltimes M:
		d\in I,\ m+N=\varphi(d)\}.
		\]
		Conversely, every triple $(I,N,\varphi)$ satisfying
		\[
		I\lhd D,\qquad N\leq_D M,\qquad IM\subseteq N,
		\qquad
		\varphi\in\Hom_D(I,M/N)
		\]
		determines an ideal
		\[
		J_\varphi
		=
		\{(d,m)\in D\ltimes M:
		d\in I,\ m+N=\varphi(d)\}.
		\]
		Furthermore,
		\[
		\pi(J_\varphi)=I.
		\]
		Thus the ideals of $D\ltimes M$ are in bijection with the triples
		$(I,N,\varphi)$ satisfying the conditions above.
	\end{theorem}

	\begin{proof}
		Let $J$ be an ideal of $R$, and let
		\[
		I=\pi(J).
		\]
		Since $\pi$ is a ring homomorphism, $I$ is an ideal of $D$.
		
		Next define
		\[
		N=\{m\in M:(0,m)\in J\}.
		\]
		It is immediate that $N$ is an additive subgroup of $M$. If
		$a\in D$ and $m\in N$, then
		\[
		(a,0)(0,m)=(0,am)\in J,
		\]
		so $am\in N$. Hence $N$ is a $D$-submodule of $M$.
		
		We claim that
		\[
		IM\subseteq N.
		\]
		Let $d\in I$ and $x\in M$. Since $d\in\pi(J)$, there is some
		$m\in M$ such that $(d,m)\in J$. Therefore
		\[
		(0,x)(d,m)=(0,dx)\in J.
		\]
		Hence $dx\in N$, proving that $IM\subseteq N$.
		
		For each $d\in I$, choose $m\in M$ such that $(d,m)\in J$, and
		define
		\[
		\varphi(d)=m+N\in M/N.
		\]
		This is well-defined. Indeed, if both $(d,m)$ and $(d,m')$ belong to
		$J$, then
		\[
		(0,m-m')=(d,m)-(d,m')\in J,
		\]
		and consequently $m-m'\in N$.
		
		The map $\varphi$ is $D$-linear. For $a\in D$ and $d\in I$, if
		$(d,m)\in J$, then
		\[
		(a,0)(d,m)=(ad,am)\in J.
		\]
		Thus
		\[
		\varphi(ad)=am+N=a(m+N)=a\varphi(d).
		\]
		
		We now prove that
		\[
		J
		=
		\{(d,m)\in D\ltimes M:
		d\in I,\ m+N=\varphi(d)\}.
		\]
		If $(d,m)\in J$, then $d\in I$ and, by definition,
		\[
		m+N=\varphi(d).
		\]
		Conversely, suppose $d\in I$ and
		\[
		m+N=\varphi(d).
		\]
		By the definition of $\varphi$, there is some $m_0\in M$ such that
		\[
		(d,m_0)\in J
		\qquad\text{and}\qquad
		m_0+N=\varphi(d).
		\]
		Hence $m-m_0\in N$, so
		\[
		(0,m-m_0)\in J.
		\]
		Therefore
		\[
		(d,m)=(d,m_0)+(0,m-m_0)\in J.
		\]
		
		For the converse construction, suppose that $I$, $N$, and $\varphi$
		satisfy the stated conditions. Let
		\[
		J_\varphi
		=
		\{(d,m):d\in I,\ m+N=\varphi(d)\}.
		\]
		It is clearly an additive subgroup of $D\ltimes M$. To prove that it
		is an ideal, let
		\[
		(d,m)\in J_\varphi
		\qquad\text{and}\qquad
		(a,x)\in D\ltimes M.
		\]
		Then
		\[
		(d,m)(a,x)=(ad,am+dx).
		\]
		Since $I$ is an ideal, $ad\in I$. Moreover,
		\[
		dx\in IM\subseteq N.
		\]
		Consequently,
		\[
		am+dx+N=am+N.
		\]
		Since $\varphi$ is $D$-linear,
		\[
		\varphi(ad)=a\varphi(d)=a(m+N)=am+N.
		\]
		Here the first component in the product is $ad$, while the second
		component is $am+dx$ when the factors are written in the order
		$(a,x)(d,m)$. Thus, using commutativity of $D$, we obtain
		\[
		am+dx+N=\varphi(ad),
		\]
		and hence
		\[
		(a,x)(d,m)\in J_\varphi.
		\]
		Therefore $J_\varphi$ is an ideal.
		
		Finally, for every $d\in I$, the definition of $J_\varphi$ provides
		an element of $J_\varphi$ whose first component is $d$. Hence
		\[
		\pi(J_\varphi)=I.
		\]
		The uniqueness of $\varphi$ follows immediately from the equality
		\[
		\varphi(d)=m+N
		\]
		for any $(d,m)\in J$. This completes the proof.
	\end{proof}

\end{document}
